\subsection{Seguimiento y Control de Fallas}

El seguimiento y control de todas las fallas reportadas se gestionará de manera centralizada para asegurar una resolución transparente y eficiente.

\subsubsection{Sistema de Tickets}
Se utilizará un sistema de gestión de tickets (por ejemplo, OTRS, Jira Service Management o similar) como herramienta principal para el registro, seguimiento y control de todas las incidencias. Cada vez que se reporte una falla, se creará un ticket único con un número de identificación. Este ticket contendrá toda la información relevante del incidente, incluyendo:
\begin{itemize}
    \item Descripción del problema.
    \item Fecha y hora del reporte.
    \item Usuario que reporta.
    \item Nivel de prioridad asignado.
    \item Técnico asignado.
    \item Todas las acciones y comunicaciones realizadas.
    \item Solución final y fecha de cierre.
\end{itemize}

\subsubsection{Comunicación y Notificaciones}
El usuario que reportó la falla recibirá notificaciones automáticas por correo electrónico en cada etapa clave del proceso:
\begin{itemize}
    \item Creación del ticket.
    \item Asignación de un técnico.
    \item Actualizaciones importantes sobre el progreso.
    \item Cierre del ticket.
\end{itemize}
Además, los usuarios podrán consultar el estado de sus tickets en cualquier momento a través del portal de soporte en línea, utilizando el número de identificación de su ticket.

\subsubsection{Monitoreo del Cumplimiento}
El administrador de la red revisará periódicamente los tickets abiertos y cerrados para supervisar el cumplimiento de los tiempos de respuesta y solución establecidos en este SLA. Los resultados de este monitoreo se incluirán en los reportes mensuales de soporte técnico.
